% =========================================================
%  CODIGO MATHEMATICA  -  anadir al preambulo SOLO si el informe lleva codigo Mathematica
%  Pegar justo antes de \begin{document}. Uso en el cuerpo:
%     \lstset{style=mathematicaStyle}
%     \begin{lstlisting} ... \end{lstlisting}
% =========================================================
\usepackage{listings}
\definecolor{codegray}{rgb}{0.5,0.5,0.5}
\definecolor{codepurple}{rgb}{0.58,0,0.82}
\definecolor{backcolour}{rgb}{0.96,0.96,0.96}
\definecolor{mma-comment}{RGB}{128, 128, 128}
\definecolor{mma-function}{RGB}{0, 0, 255}

\lstdefinestyle{mathematicaStyle}{
    language=Mathematica,
    backgroundcolor=\color{backcolour},
    commentstyle=\color{mma-comment}\itshape,
    keywordstyle=\color{mma-function}\bfseries,
    emph={Plot3D, Do, AppendTo, Table, Graphics3D, Line, Show, Module, Evaluate,
          ClearAll, Map, D, Legended, SwatchLegend, Thickness, PointSize, Point,
          Axes, AxesLabel, Boxed, ViewPoint, ImageSize, PlotLabel, Style},
    emphstyle=\color{mma-function}\bfseries,
    stringstyle=\color{codepurple},
    basicstyle=\ttfamily\footnotesize,
    breaklines=true,
    numbers=left,
    numbersep=5pt,
    numberstyle=\tiny\color{codegray},
    frame=single,
    rulecolor=\color{black!30},
    morecomment=[s]{(*}{*)},
    showstringspaces=false
}
